\begin{explanationbox}
	\n\n\textbf{\textcolor{CExplanation}{一句话直觉}}\n\n动点到两个固定点距离之比为常数时，其轨迹是一个圆；比值越大或越小，圆的位置和大小随之改变.\n\n

	\medskip
	\n\textbf{\textcolor{CExplanation}{核心拆解}}\n
	\begin{description}[style=nextline, leftmargin=2.8em, labelsep=0.5em, itemsep=0.45em]
		\n
		\item[\textbf{要点一（内分外分点）：}]
			直径的两个端点分别是线段 $AB$ 按比例 $k$ 的内分点 $C$ 和外分点 $D$.\n
		\item[\textbf{要点二（圆心半径公式）：}]
			圆心在直线 $AB$ 上，半径 $R = \dfrac{k}{|1-k^{2}|}\,|AB|$，圆心到 $A$ 的有向距离为 $\dfrac{k^{2}}{k^{2}-1}\,|AB|$.\n
	\end{description}
	\n\n

	\medskip
	\n\textbf{\textcolor{CExplanation}{几何本质（定比分圆）}}\n\n阿波罗尼斯圆是通过定比分点构造的圆：以 $AB$ 的内、外分点 $C,D$ 为直径端点，则圆上任一点 $P$ 都满足 $PA:PB=k$.\n\n

	\medskip
	\n\textbf{\textcolor{CExplanation}{代数意义（距离平方转化）}}\n\n设 $A(0,0),B(d,0)$，由 $PA^{2}=k^{2}PB^{2}$ 整理可得圆的标准方程，从而得到圆心和半径.\n\n

	\medskip
	\n\textbf{\textcolor{CExplanation}{考点价值（轨迹与最值）}}\n
	\begin{itemize}[leftmargin=*, itemsep=0.5em, topsep=0.4em]
		\n
		\item \textbf{考法一（求轨迹方程）：} 已知两点和距离比，直接套用结论得圆方程.\n
		\item \textbf{考法二（三角形最值）：} 利用轨迹圆求三角形面积或线段和最值.\n
	\end{itemize}
	\n\n

	\medskip
	\n\textbf{\textcolor{CExplanation}{顿悟点}}\n\n关键是将距离比转化为距离平方等式，通过代数变形自然导出圆的方程；同时要重视 $k=1$ 的退化情形.\n\n

	\medskip
	\n\textbf{\textcolor{CExplanation}{使用场景}}\n
	\begin{itemize}[leftmargin=*, itemsep=0.5em, topsep=0.4em]
		\n
		\item \textbf{场景一（求轨迹方程）：} 已知两点和距离比，直接套用结论得出圆的方程.\n
		\item \textbf{场景二（解三角形最值）：} 利用阿波罗尼斯圆将边比条件转化为圆的轨迹，再求最值.\n
	\end{itemize}
	\n
\end{explanationbox}
